\subsubsection{The nonidentity terms with vanishing single-term Gram invariant}
\label{subsec:zero-gram-diagonal}

Fix one word \(u\). We first sum the terms of \(K''_{uu}\) on the actual marked row. Two different ranks occur in this calculation. The integer \(j\) below belongs to a single pair of flags defining an orbital term. The corank of a pair of orbital terms enters later, after the terms with zero Gram invariant have been removed.

\begin{definition}[The minor label of an orbital term]\label{def:zero-gram-label}
Put \(s=m/2-1\), \(d=\operatorname{rd}(E)\), and write \(v\) for one complex embedding above each real place of the maximal real subfield. In the two ordered isotropic flags write \(Z_h=(\langle e_a,e'_b\rangle)_{a,b\le h}\). Let \(j=\operatorname{rank}_E Z_s\), \(\delta_0=1\), \(\delta_h=\det Z_h\) for \(1\le h\le j\), and \(\delta_h=0\) for \(j<h\le s\). For the resulting label \(\delta\) set
\[
 e_h=\sum_{a=1}^h(s+1-a),\quad
 a_{v,h}=|v(\delta_h)|d^{-2e_h},\quad
 D_\delta=\prod_v a_{v,j}^2.
\]
For \(H_1=b^{C_Hm}\), with \(C_H\) the common length constant below, put
\[
 S_v=\max\left(1,\max_{h\le j}\frac{a_{v,h}}{(2H_1)^{2s}}\right),
 \quad S_\delta=\prod_v S_v,\quad H_C=b^{Cm^2f}.
\]
\end{definition}

\begin{auxlemma}[The positive rank of a nonidentity term]\label{lem:zero-gram-positive-rank}
For a term of \(K''_{uu}\) with nonzero coefficient, the isotropic flags have a common tail after unipotent changes of frame. The prefix pairing consists of invertible scalar multiples of reversed identity matrices on consecutive blocks. In particular
\[
                 1\le j\le s,\qquad \delta_j\ne0.
\]
\end{auxlemma}
\begin{proof}
The finite integral restricts to supported flags. Translating the real integral by a member of their common unipotent changes it by the ratio of the two characters. A nonzero coefficient therefore requires their equality on this intersection. Here the last vector in each augmented flag has negative norm \(-T_u\). 

Write \(X_a=\langle e_1,\ldots,e_a\rangle\), \(X=X_s\), and similarly for the other flag, with \(e_{s+1}=z\), \(e'_{s+1}=z'\). Put \(I=X\cap X'\). The complement of the distinguished negative line in \(X^\perp/X\) is positive definite. Thus
\[
   y\in(X+Ez)^\perp,\quad\langle y,y\rangle=0
                      \quad\Longrightarrow\quad y\in X .
\]
Also \(z\notin X+X'\): if \(z=x+x'\) with \(x\in X,x'\in X'\), orthogonality of \(z\) to \(x\) gives \(\langle x',x\rangle=0\), contradicting \(\langle z,z\rangle=-T_u\). Both facts hold on exchanging the flags.

Suppose first that \(0\ne w\in I\) has leading terms \(\alpha e_a\), \(\alpha'e'_{a'}\) in the two filtrations. For \(v\perp X_a+X'_{a'}\), the infinitesimal shear
\[
       Y(\xi)=\langle\xi,v\rangle w-\langle\xi,w\rangle v
\]
belongs to both unipotent Lie algebras. Its character derivatives read the coefficients \(\alpha\langle e_{a+1},v\rangle\) and \(\alpha'\langle e'_{a'+1},v\rangle\). All earlier superdiagonal entries vanish by the two orthogonality conditions. Varying \(v\) through its scalar multiples and using the nondegenerate trace pairing gives
\[
  \alpha e_{a+1}-\alpha'e'_{a'+1}\in
    ((X_a+X'_{a'})^\perp)^\perp=X_a+X'_{a'} .
\]


As long as both successors are isotropic, moving the two lower terms to their respective sides constructs a vector in \(I\) with leading coefficients \(\alpha,\alpha'\) on the next steps. If \(a\ne a'\), one side reaches its negative vector first, placing it in \(X+X'\). Thus \(a=a'\). At the end
\[
       z'-cz\in X+X',\qquad c=\alpha/\alpha' .
\]
This quotient determines \(c\) uniquely since \(z\notin X+X'\). Starting from the least common leading index \(j+1\) constructs every subsequent index. Subtracting these vectors in descending order proves \(X=X_j\oplus I\), \(X'=X'_j\oplus I\).

If \(x\in X_j\) pairs to zero with \(X'_j\), then \(x\perp X'\) and \(x\perp z'\), the latter by \(z'-cz\in X+X'\). The positivity implication for the primed flag gives \(x\in I\), hence \(x=0\). Therefore the prefix pairing is invertible.

For \(I=0\), take \(j=s\). Write \(d_a=\det Z_a\), now including \(z,z'\) as the last entries at \(a=s+1\). The single-term condition in the chosen normalization is
\[
      |d_{s+1}|^2-T_u^2|d_s|^2=0 .
\]


If \(d_s=0\), a nonzero \(x\in X\) annihilates \(X'\). Necessarily \(\langle x,z'\rangle\ne0\) by positivity and \(I=0\). But \(d_{s+1}=0\) gives a nonzero \(x'+\beta z'\), \(x'\in X'\), orthogonal to \(X+Ez\). Here \(\beta\ne0\) by the same positivity argument with the frames interchanged. These relations imply the contradiction
\[
        0=\langle x,x'+\beta z'\rangle
             =\bar\beta\langle x,z'\rangle\ne0 .
\]
Therefore \(d_s\ne0\).

 Remove the projections of \(z,z'\) onto the hyperbolic space \(X+X'\); these projections lie in \(X,X'\), respectively, so their removal preserves both negative norms. The Schur complement has absolute value \(|d_{s+1}/d_s|=T_u\). Equality in the reversed Cauchy--Schwarz inequality for the two resulting negative vectors makes them proportional: subtract the projection onto the first negative line; the residual norm in its positive definite perpendicular is zero. This again gives \(z'-cz\in X+X'\).

On the paired prefix, write the pivot in column \(b\) as \(t_b\) in row \(\pi(b)\). Conjugate an elementary shear between columns \(c<b\) with \(\pi(c)=\pi(b)+1\) into the opposite frame. Equality of characters for every scalar gives
\[
              b=c+1,\qquad t_c=t_b .
\]
Starting in the first column, locate the preceding pivot row, and repeat until row one is reached. This produces the first reversed block. Deleting it gives the same rule in the remaining columns. On this first block, of size \(h=\pi(1)\), the pivot relations read
\[
      \pi(b)=h+1-b,\qquad t_b=t_1\quad (1\le b\le h).
\]

In both cases the unitriangular changes and the clearing of the active pairings place the two terminal vectors on the same negative line. For \(I\ne0\) one first takes \(I^\perp/I\) and splits off the hyperbolic prefix. Its projections of \(z,z'\) lie on the respective isotropic sides, since each vector is orthogonal to its own flag. Removal preserves the norms and leaves vectors differing by the scalar \(c\). Hence
\[
               -T_u=-T_u c\bar c .
\]

If \(j=0\) all isotropic flag steps and both enlarged final spaces coincide. Primitivity of the graded lattices gives \(c\in\mathcal O_E^\times\). Every conjugate has modulus one. The characteristic polynomials of the powers of \(c\) have bounded integral coefficients, so \(c\) has finite order. 

The markings imply
\[
     c\equiv1\pmod {\mathfrak g_i}\quad(\mathfrak g_i\mid g_i,\ i=0,1),
     \qquad (T_u,g_0g_1)=1 .
\] The reduction kernel on roots of unity at \(g_i\) has \(g_i\)-power order: for an order prime to \(g_i\), the factor \((c^a-1)/(c-1)\) is a unit when \(c\equiv1\) and \(g_i\nmid a\). Thus \(c=1\).

For equal graded vectors the two integral lifts at each step differ in the preceding flag lattice. Successive upper triangular shears remove these differences. The graded ideals are projective over \(\mathcal O_E\), so each shear lifts from the graded quotient; the dual action is the inverse adjoint. The residual difference of distinguished vectors lies in the common isotropic space, as allowed in the identity class. This is precisely the identity part removed in Definition~\ref{hyp:full-gram-row}.
\end{proof}

\begin{auxlemma}[Decay at fixed minor label]\label{lem:zero-gram-fixed-label}
Write \(w_\Lambda=1/|\Delta_\Lambda|\) for the definite row mass and divide the real coefficient vector by its endpoint factor \(\sqrt{w_\Lambda w_{\Lambda'}}\), where \(\Lambda,\Lambda'\) are its two endpoint rows. Denote its resulting Hilbert--Schmidt norm by \(A_t\). For an adapted rational normal form let \(\Gamma_0\) be its reference integral stabilizer and \(\Gamma_\Lambda\) the actual integral stabilizer of the term. With \(J=\Gamma_0\cap\Gamma_\Lambda\), put
\[
 \omega_t=\inf_{\text{adapted forms}}
          \frac{[\Gamma_0:J]}{[\Gamma_\Lambda:J]} .
\]
Let \(C_p=16\). There are fixed \(C_H,C\) and a scale threshold, independent of the retained row and word, for which
\begin{equation}\label{eq:zero-gram-angular-gain}
 \frac{A_t}{\sqrt{\dim V_u}}\le
 \omega_t (C/k)^{j(k-1-j)f}
       D_\delta^{-(k-1-j)/2} H_C^{Cs^4}
                    S_\delta^{-(C_p+4)k}.
\end{equation}
\end{auxlemma}

\begin{proof}
\proofstep{Moments on each side of the integral.}
For the orthogonalized lengths \(l_0,\ldots,l_s\) of either frame set
\[
 R(x)=\max(1,l_0,l_0^{-1},\ldots,l_s,l_s^{-1}),\qquad
 q=\lceil C_1ks\rceil ,
\]
where \(C_1\) is a sufficiently large fixed numerical constant. Use the lowest type with initial shift zero in the first-length and successive-ratio recursions in Proposition~\ref{prop:paired-frame-moments}. Their iteration is valid for all sufficiently large selected ranks since
\[
       (s+1)q\le\min(km^5,2(k+m-1)^2).
\]

 For inverse lengths the terminal vector has negative norm: subtracting vectors of its isotropic flag preserves this norm. The Euclidean majorant of the squared norm is consequently at least its absolute value. Thus \(l_s\ge \sqrt{T_u}\) in the original frame, and transport to the standard character costs at most \(b^{C m}\). Use
\[
 l_a=l_0\prod_{i<a}(l_{i+1}/l_i),\qquad
 l_a^{-1}=l_s^{-1}\prod_{i=a}^{s-1}(l_{i+1}/l_i).
\]
H\"older with at most \(s+1\) equal exponents gives a common \(H_1=b^{C_Hm}\) in the two \(2q\)-moments of every length. Increasing \(C_H\) absorbs \(2s+3\) on taking the \(2q\)-th root. If \(\mu_F=|F|^2d\nu/\int|F|^2d\nu\), Markov's inequality gives
\[
       \mu_F\{R>2^a H_1\}\le 2^{-2aq}\qquad(a\ge0).
\] Use the trivial bound on \(R(x)\le 2H_1\).

\proofstep{Volume in the noncommon directions.}
Work at a single real place in an adapted Witt decomposition with positive majorant. Write \(D=a_{v,j}\), \(D_e=\max(1,\max_{a_{v,h}>0}a_{v,h}^{-1})\), and \(n=k-1-j\). The two frames have \(j\) noncommon unipotent columns \(B_1,\ldots,B_j\in\mathbb C^{k+1}\). 

In Iwasawa coordinates remove the last coordinate and the direction of the transported negative vector by orthogonal projection onto \(P\). These impose at most two complex linear conditions, so \(p=\dim P\ge k-1\). In the transformed middle coordinates put
\[
 Y=\|P B_1\wedge\cdots\wedge P B_j\| .
\]
Equivalently, the projected wedge of the first \(j\) opposite frame vectors has norm \(DY\): the normal-form change of its \(j\) columns has determinant of modulus \(D\). Define \(Y'\) by exchanging the frames. If both length heights are at most \(R_0\), put
\[
                      M=R_0^{2s}D_e .
\]

 Solving for a unipotent column by the prefix pairing uses a cofactor divided by its block-end minor. The numerator is bounded by the products of the two prefix volumes, at most \(R_0^{2s}\), by the exterior-power Cauchy--Schwarz inequality. In the transverse coordinates the resulting boxes have radii
\[
   \|B_i\|\le M,\qquad |u_{ab}|\le M,\qquad
   |v_{ab}|,\ |\Im v_{aa}|\le 2sM^2 .
\]
Here \(u\) is the upper triangular change between blocks, and \(v\) is the central skew-Hermitian coordinate. Its quadratic correction is half the pairwise Gram matrix of the columns. Triangular coordinates have unit Jacobian, and only \(O(s^2)\) complex or real entries occur outside the \(B_i\).



To bound the volume where \(0<Y\le 2y\), integrate the columns in order. Project the \(i\)-th column perpendicular to the previous \(i-1\) projected columns. Writing \(p_i=p-i+1\), its length \(r_i\) has radial measure
\[
             \frac{2\pi^{p_i}}{(p_i-1)!}r_i^{2p_i-1}dr_i,
             \qquad p_i>n.
\]
The exceptional dependent columns have measure zero. Since \(p-i+1>n\), integration of \(r_i^{-2n}\) on the ball of radius \(M\) is finite. On \(\prod_i r_i\le2y\) insert \((2y/\prod_i r_i)^{2n}\). 

The remaining powers of \(M\) from each full column are \(2(k+1-n)=2j+4\). The disk integrals in the at most \(i+1\) removed complex directions cost at most \(C^{i+1}\) apart from these powers of \(M\). Using \(n!\ge(n/e)^n\) gives the transverse volume bound
\begin{equation}\label{eq:zero-gram-fibre}
 (Cs)^{Cs^2}M^{Cs^2}(C/n)^{jn}(2y)^{2n}.
\end{equation}

\proofstep{Pairing the two transverse volumes.}
Write \(Z\) for the nonsingular \(j\)-square prefix pairing and \(U,U'\) for the matrices of the two opposite projected frames before division of their wedge norms by \(D\). Up to the choice of Hermitian convention the mixed form on the two positive projections has matrix
\[
                         U (Z^{-1})^{\mathsf t} U'^* .
\]
It is a difference of two pairings. To see this, split a projected vector into its component orthogonal to the isotropic flag and its isotropic component. In the latter term solve for the prefix coefficients using \(Z\); the common tail pairs to zero. 

The two remaining pairings are the original positive-majorant pairing and the pairing on the positive definite complement of the negative line. The latter components have norm at most that of the test vectors: the discarded isotropic component is orthogonal to the definite component for the Hermitian form and has norm zero for that form. Writing \(B_1,B_2\) for these two forms, Cauchy--Schwarz gives
\[
 |B_i(\xi,\eta)|\le\|\xi\|\,\|\eta\|,\qquad
                  \|B_1-B_2\|_{\rm op}\le2 .
\] Taking the \(j\)-th exterior power therefore gives
\[
 \frac{\|\wedge^j U\|\,\|\wedge^j U'\|}{|\det Z|}
                 =DYY'\le 2^j .
\]


Cauchy--Schwarz on the unfolded integral uses the square root of each of the two bounds \eqref{eq:zero-gram-fibre}. More explicitly, integrate first along the stabilizer of each individual frame. The Bessel norm is constant on this fibre and the quotient by the common stabilizer has the transverse coordinates above. Applying Fubini on both sides of Cauchy--Schwarz leaves the full squared Bessel norm on each individual quotient. The same integer \(n=k-1-j\) occurs in both volumes.

For fixed dyadic length heights let \(R_0\) be the larger endpoint. Then \(Y,Y'\le M^j\) by Hadamard. On dyadic wedge shells with lower endpoints \(y,y'\), sum \((yy')^n\) subject to \(yy'\le 2^j/D\). 

For fixed product the number of shells is at most \(1+\log_2^+(M^{2j}/(yy'))\); summation in successively halved products gives
\[
       \sum (yy')^n\le C^{jn}D^{-n}
                (1+\log^+(DM^{2j})).
\]
The logarithm costs at most an extra polynomial factor in \(s,R_0,D_e\): the nonzero minor \(D\) is bounded by \(R_0^{2s}\). 

The dyadic subdivisions therefore give a factor \((C/k)^{jn}D^{-n}\) per place, since \(n\ge k/2\). On the length shells with indices \(a,a'\) the remaining multiplier is at most
\[
              (CsH_1)^{Cs^3}D_e^{Cs^2}2^{Cs^3(a+a')}.
\]

Hadamard applied separately to the two prefixes also gives \(S_v\le 2^{s(a+a')}\). Use a portion of the moment decay \(2^{-q(a+a')}\) to extract \(S_v^{-(C_p+6)k}\). The rest sums geometrically, since \(k\gg s^4\) and \(C_1\) can be chosen after all the absolute box constants. 

For every nonzero label minor the graded ideals of the two flags give \(\delta_h\in\mathfrak d^{2e_h}\), where \(\mathfrak d\) denotes the different. Hence \(\prod_v a_{v,h}\ge1\). In particular
\[
 \prod_v D_{e,v}\le\prod_{\delta_h\ne0}\prod_v\max(1,a_{v,h})
                 \le ((2H_1)^{2sf} S_\delta)^{s+1}.
\]
Indeed for a nonzero minor
\[
       \prod_v\max(1,a^{-1}_{v,h})
             =\frac{\prod_v\max(1,a_{v,h})}{\prod_v a_{v,h}} .
\] Its contribution to the exponent of \(S_\delta\) is \(O(s^3)\), absorbed by the two spare multiples of \(k\). Multiplication of \(D^{-n}\) over the places gives \(D_\delta^{-n/2}\).

\proofstep{The common stabilizer and quotient measures.}
Write \(v_E=d/2\). The middle directions of the reference stabilizer are the \(s-j\) common columns with coefficients in the rank-\(k\) definite lattice. In the scaled graded coordinates each such column has covolume \(v_E^{kf}\): the lattice is unimodular and the powers of the different in its graded ideal cancel the frame dilation. In the ordinary complex coordinates the scalar formula underlying this computation is
\[
      \operatorname{covol}(\mathcal O_E)=2^{-f}|D_E|^{1/2}=v_E^f .
\]

 The upper and skew-Hermitian layers have at most \(2(s+1)^2\) real coordinates per place. An off-diagonal complex entry is in its graded fractional ideal; a diagonal skew entry is in the trace-zero part. 

Taking successive quotients by the central entries eliminates the quadratic cross terms. Clearing a factor of two on each skew diagonal suffices for the half-Gram correction. Pairing conjugate off-diagonals and using the dual graded ideal bounds the remaining factors by \((8v_E)^{2(s+1)^2 f}\) for the common stabilizer, with one further such factor for the two frame normalizations together. The resulting bound, including the Jacobian from both real frames, is
\[
  \omega_t v_E^{k(s-j)f}(8v_E)^{4(s+1)^2f}
  \left((2v_E)^{-sk}2^{s(k+1)+s(s-1)}\right)^f .
\]


Indeed, comparison on the common finite-index subgroup \(J\) multiplies the reference volume by \([\Gamma_0:J]/[\Gamma_\Lambda:J]\); the bound is valid in every adapted frame. The high powers of \(v_E\) and two in the last display simplify to \(v_E^{-kjf}2^{s^2f}\le2^{s^2f}\). Since \(2\le d\le b\), all remaining height costs, including \((CsH_1)^{Cs^4f}\), are at most \(H_C^{Cs^4}\) for a fixed larger \(C\). The two squared Bessel integrals divided by the identity normalization have product one for the operator norm; passing to Hilbert--Schmidt norm costs \(\sqrt{\dim V_u}\). Taking the infimum in the index ratio proves \eqref{eq:zero-gram-angular-gain}.
\end{proof}

For the definite count, write \(R=\mathcal O_E\) and
let \(W\) denote the definite Hermitian space of rank \(k\) obtained
from the common target row. For a full integral lattice \(A\) in a
nondegenerate subspace, write \(A^\#\) for its Hermitian dual
in that subspace and \(\Delta(A)=|A^\#/A|\). In particular,
\(\Delta(0)=1\). Fix the genus \(\mathcal G\) of
Proposition~\ref{prop:flag-row-mass}, with \(T=T_u\). Thus
\[
     L^\#/L\simeq R/TR,\qquad D=\Delta(L)=T^{2f}
          \qquad(L\in\mathcal G),\qquad f=[F:\mathbf Q].
\]

An \emph{embedded datum of index \(x\)} over \(L\) consists of
a full integral lattice \(B\subset W\) of discriminant order \(D\)
and the inverse image \(\Gamma\subset B\times L\) of the graph
of a bijective isometry \(B/gB\to L/gL\) of the forms modulo \(g\).
Its diagonal lattice must satisfy
\[
 M_\Gamma=\{z\in W:(z,z)\in\Gamma\},\qquad
 [L:M_\Gamma]=x.
\]

The quotient \(L/M_\Gamma\) must be generated by at most \(s=m/2-1\)
elements over \(R\).
Write \(\mathcal T(L;x,x')\) for the set of triples consisting of
two such data \(d=(B,\Gamma)\), \(d'=(B',\Gamma')\), of indices
\(x,x'\), and a rational unitary isometry \(\eta:W\to W\) such that
\[
          \eta B'=B,\qquad (\eta\times1)\Gamma'=\Gamma .
\]

The lattices and isometries here are counted individually inside
the one target space. Normalization by target automorphisms occurs
only when averaging over the genus.

\begin{auxlemma}[Compatible embedded data in the definite genus]
\label{lem:zero-gram-definite-mean}
Let \(k\ge3\), \(g\ge2\), and \(x,x'\ge1\) be integers, and let
\(E/F\), \(T=T_u\) and \(\mathcal G\) be as in
Proposition~\ref{prop:flag-row-mass}. Put
\[
 \begin{gathered}
   c=1+\log_2(2f),\qquad v=\mathrm{rd}(E)/2\ge1,\qquad
   X=xx',\qquad Y=XD^2,\\
   M_b=\bigl(20f(2v)^2\bigr)^{2f(6+c)},\qquad
   \mathcal A=M_b(2v)^f,\\
   U_*=(32/g)^{kf}(4kf)^{4f}Y^7,\qquad
   V_*=(2/g)^{kf}(4kf)^{4f}DX .
 \end{gathered}
\]
For the fields of \eqref{eq:host-character-compositum},
the quadratic subfield \(\mathbf Q(\sqrt{-d})\) has discriminant
of absolute value \(d\equiv7\bmod16\) by
Proposition~\ref{prop:quadratic-prime-pair}. The discriminant tower formula
gives \(2v=\mathrm{rd}(E)\ge\sqrt d\).
For the probability proportional to \(1/|\operatorname{Aut}(L)|\)
on isometry classes in \(\mathcal G\),
\begin{equation}\label{eq:zero-gram-pair-mean}
   \big\langle |\mathcal T(L;x,x')|\big\rangle_{\mathcal G}
   \le X^{k/2}Y^{6s+10+4c}
      \left(1+\mathcal A\sum_{h=1}^{k-1}U_*^h+V_*^k\right).
\end{equation}
The bound applies when the two quotients have different minimal
numbers of generators, each at most \(s\).
\end{auxlemma}

\begin{proof}
\proofstep{Data with a prescribed diagonal lattice.}
Take \(M=M_\Gamma\) and \(u_0=2s+1\). The inclusions
\[
 M\subset L\subset L^\#\subset M^\#,\qquad
 |M^\#/M|=Dx^2,\qquad |M^\#/L|=Dx
\]
give \(M^\#/M\) at most \(u_0\) generators. Indeed the two outside
quotients are dual up to conjugation and the middle quotient is cyclic.
Every submodule of this finite module has the same generator bound
by elementary divisors locally and the Chinese remainder theorem.
Since \(M\subset B\subset M^\#\), a list of \(u_0\) elements
of \(M^\#/M\) suffices to specify \(B/M\). Once \(B\) is chosen,
fix lifts \(b_i\) of a generating list for \(B/M\). If \(l_i\)
represents the image of \(b_i\) modulo \(gL\), testing the isometry
on \(b_i\) and each \(z\in M\) gives
\[
 \langle b_i-l_i,z\rangle\in gR,\qquad
             (b_i-l_i)/g\pmod L\ \in M^\#/L .
\]

These classes determine the images of the \(b_i\), and the graph
already sends each \(z\in M\) to \(z\bmod gL\). The number of
possible data above \(M\) is consequently at most
\[
        (Dx^2)^{u_0}(Dx)^{u_0}\le(Dx)^{3u_0}.
\]

\proofstep{The common fixed space and its discriminant.}
Write \(M'=M_{\Gamma'}\), \(U=\ker(1-\eta)\), \(h=k-\dim_E U\).
On \(U\) the two graph conditions coincide. Put
\[
 M_0=M\cap U=M'\cap U,\qquad a=\Delta(L\cap U),\qquad
 y=[L\cap U:M_0]\le\min(x,x').
\]

Let \(P:W\to U^\perp\) be orthogonal projection. For \(z'\in M'\)
and \(z\in M\), the pairs \((\eta z',z')\) and \((z,z)\) belong
to \(\Gamma\). Comparing their pairings in its two coordinates gives
\[
    C(PM')\subseteq(PM)^\#,\qquad
                C=(1-\eta)|_{U^\perp}/g .
\]

Definiteness makes \(C\) invertible on \(U^\perp\). To take indices
also for these projected lattices, let \(\delta(A)\) be the real
covolume in the product of the \(f\) complex Hermitian spaces,
divided by \(v^{f\operatorname{rank}_E(EA)}\). This divides out
the covolume \(2^{-f}\sqrt{|\operatorname{disc}(E)|}=v^f\)
of \(R\) in one Hermitian line. The determinant formula gives
\[
  \delta(A^\#)=\delta(A)^{-1},\qquad
  \delta(M)=\delta(M_0)\delta(PM),\qquad
  \delta(M')=\delta(M_0)\delta(PM').
\]

The last two identities follow by projecting a real fundamental
parallelepiped along \(U\). For integral \(A\),
\(\delta(A)^2=\Delta(A)\). At each complex embedding, every
eigenvalue of \(\eta\) has absolute value one. With \(q_0=2/g\)
the projected inclusion therefore yields
\begin{equation}\label{eq:zero-gram-fixed-discriminant}
  ay^2=\delta(M_0)^2
   \le |N_{E/\mathbf Q}\det C|\,\delta(M)\delta(M')
   \le q_0^{2fh}DX .
\end{equation}

In particular, only integers \(1\le a\le DX\) can occur.

\proofstep{Triple count for a specified space \(U\).}
A rank-\(n\) projective \(R\)-lattice has at most \(I^{n+c}\)
sublattices of index \(I\), by the prime-by-prime chain count in
Proposition~\ref{prop:primitive-flag-means}. Thus for a fixed \(y\)
the number of \(M_0\) is at most
\[
   y^{k-h+c}\le (q_0^{2fh}DX/a)^{k/2}X^c.
\]
There are at most \(X\) choices for \(y\). Given \(M_0\), the image
of \(M\) in \(L/(L\cap U)\) has index \(x/y\), hence at most
\((x/y)^{h+c}\) possibilities. For a fixed image \(J\), lifts form
a torsor under \(\operatorname{Hom}_R(J,(L\cap U)/M_0)\).
Locally \(J\) is free of rank \(h\), so
\[
 |\operatorname{Hom}_R(J,(L\cap U)/M_0)|=y^h,\qquad
             (x/y)^{h+c} y^h\le x^{2h+c}.
\]
This bounds the choices for \(M\). Replacing \(x\) by \(x'\) gives
the corresponding bound for \(M'\).

With both data fixed, all compatible \(\eta\) that fix \(U\)
pointwise lie in one coset of the group of automorphisms of \(B\)
fixing \(U\) pointwise. This finite group acts faithfully on
\(B/(B\cap U)\), since orthogonality identifies its rational
action with the action on \(U^\perp\). For \(h>0\), put \(n=2fh\).
By Minkowski's theorem \cite[\S1.1]{serre-finite-orders}, its
order divides \((2n)!\). Indeed the bound on the exponent at
each prime \(p\) satisfies
\[
 \sum_{i\ge0}\left\lfloor\frac{n}{p^i(p-1)}\right\rfloor
       \le \sum_{i\ge1}\left\lfloor\frac{2n}{p^i}\right\rfloor
       =\operatorname{ord}_p((2n)!) .
\]
Consequently there are at most \((4hf)^{4hf}\) isometries, with
the value \(1\) when \(h=0\). Multiplying the bounds for the
intersection, its two lifts, the two graphs and the isometry gives
the following count over a specified \(U\):
\begin{equation}\label{eq:zero-gram-one-space}
  (q_0^{2fh}DX/a)^{k/2}
        X^{2h+2c+1}Y^{3u_0}(4hf)^{4hf}.
\end{equation}

\proofstep{The mean number of fixed spaces in intermediate dimension.}
For \(1\le h<k\), write \(z=k-h\), \(H_0=aD\) and
\(\lambda=\log_2(2f)=c-1\). Such spaces \(U\) of discriminant order
\(a\) correspond bijectively to their primitive lattices \(L\cap U\).
We need the constant in the gluing calculation of
Proposition~\ref{prop:primitive-flag-means}. Write \(m_i\) for
the standard unimodular mass in positive rank \(i\). The complementary
lattice \(L\cap U^\perp\) has discriminant order \(b\le H_0\):
the gluing group injects into both discriminant modules. In fact,
\[
    H=L/((L\cap U)\oplus(L\cap U^\perp)),\qquad
    ab=D|H|^2,\qquad |H|\le a .
\]

For the rational isometry type of \(U\), signatures are fixed and
at the finite places the determinant modulo norms classifies the
Hermitian space \cite{landherr}. At split places the class is trivial. At unramified
inert places it is determined by valuation parity and is trivial
away from primes over \(a\). At an inert prime \(\mathfrak P\)
an integral basis with Gram matrix \(H_{\mathfrak P}\) gives
\[
 \operatorname{ord}_{\mathfrak P}\det H_{\mathfrak P}
 =\operatorname{length}_{R_{\mathfrak P}}
     ((L\cap U)^\#/(L\cap U))_{\mathfrak P}.
\]

Thus these lengths also specify the determinant parities. The prime lists for those modules, of
total order \(a\), number at most \(a^\lambda\) by the index
count in Proposition~\ref{prop:primitive-flag-means}.
Local-to-global classification and Witt cancellation determine the
complementary rational space for each choice. Allow \(a^{1+\lambda}\)
choices. For each type the two fixed-discriminant mass bounds and
the sum over \(1\le b\le H_0\) cost at most
\[
          m_z a^{(z+\lambda)/2},\qquad
          H_0 m_h H_0^{(h+\lambda)/2}.
\]

To glue the two lattices, choose a subgroup of the direct sum of
their discriminant modules which injects into each factor. Its
projection into the rank-\(h\) discriminant module requires at most
\(h\) generators locally, hence globally. Even allowing
\((aH_0)^{2h+1}\) choices of subgroup suffices. Orbit--stabilizer
divides this count by the two automorphism orders before summing.
With \(N_{h,a}(L)=\#\{U:\dim U=z,\ \Delta(L\cap U)=a\}\)
the resulting unnormalized mass is at most
\[
 m_zm_h\,a^{z/2}H_0^{(9h+8+4\lambda)/2}.
\]

This calculation uses positive ranks \(z,h\), so includes \(h=1\)
and \(z=1\). By Proposition~\ref{prop:flag-row-mass},
the normalizing mass is
\[
     m_k\big((T^k-1)/(T-1)\big)^f\ge m_kD^{(k-1)/2}.
\]

Since \(a^{z/2} D^{-(k-1)/2}=(a/D)^{k/2}D^{1/2}a^{-h/2}\),
division gives
\[
 \langle N_{h,a}\rangle_{\mathcal G}
   \le \frac{m_zm_h}{m_k}(a/D)^{k/2}(aD)^{5h+5+2c}.
\]

It remains to control the rank dependence of the mass ratio. Put \(r=\min(z,h)\), \(t=k-r\) and \(R_i=m_i/m_{i-1}\)
for \(i\ge2\). Shimura's adjacent-rank formula as used in
Proposition~\ref{prop:primitive-flag-means} gives
\[
 R_i=\big[(i-1)!(v/\pi)^i\big]^f(2v)^{-f/2}
        L_F(i,\chi_{E/F}^i),\qquad
 2^{-f}\le L_F(i,\chi_{E/F}^i)\le2^f .
\]
The bounds follow from the Euler product and
\(\zeta(i)^f\le\zeta(2)^f<2^f\). The rank-one estimate there,
\(m_1\le2^fh_E\), and its Minkowski ideal count
\(h_E\le(2v)^{f(1+\lambda)}\) give \(m_1\le M_b\).
Using \(v\ge1\), factorial monotonicity, and \(r+t=k\),
\[
 \begin{aligned}
  \frac{m_zm_h}{m_k}
     &=\frac{m_1}{R_{t+1}}\prod_{i=2}^r\frac{R_i}{R_{t+i}}\\
     &\le m_1(2v)^f4^{rf}\pi^{(rt+1)f}
       \le \mathcal A\,16^{hkf}.
 \end{aligned}
\]
Indeed \(r(t+1)+1\le2hk\) and \(\pi<4\).

\proofstep{Summing the three ranges.}
For \(1\le h<k\), multiply the mean just obtained by
\eqref{eq:zero-gram-one-space}. The discriminant factors cancel
in the main exponent:
\[
  (q_0^{2fh}DX/a)^{k/2}(a/D)^{k/2}
         =q_0^{fhk}X^{k/2}.
\]
Use \(X\le Y\), \(aD\le Y\) and \(3u_0=6s+3\).
For each \(a\) the result is at most
\(X^{k/2}Y^{6s+9+4c}\mathcal A U_*^h\).
There are at most \(DX\le Y\) possible orders \(a\).

When \(h=0\), \(U=W\) and \(a=D\). Direct substitution in
\eqref{eq:zero-gram-one-space} gives at most
\(X^{k/2}Y^{6s+10+4c}\).
When \(h=k\), \(U=0\) is the sole space. Counting the two
diagonal lattices before imposing compatibility costs at most
\(X^{k+c}\). With the data and isometry bounds this gives
\[
       X^{k+c}Y^{3u_0}(4kf)^{4kf}.
\]
If any such triple occurs, \eqref{eq:zero-gram-fixed-discriminant}
with \(a=y=1\) gives \(1\le q_0^{2kf}DX\), hence
\(1\le X^{k/2}(q_0^{kf}D)^k\). Multiply by the latter factor
and collect powers to obtain at most
\(X^{k/2}Y^{6s+10+4c}V_*^k\). Adding the two endpoints and
the intermediate sum proves \eqref{eq:zero-gram-pair-mean}.
\end{proof}


We now attach the embedded data to actual pairs of flag terms.
Write \(s=m/2-1,\ d=\mathrm{rd}(E)\), and let \(\mathfrak d\)
be the different used for the integral frames. Fix a word \(u\).
For a nonunit term \(t\) with \(G_t=0\) and nonzero coefficient, let \(j\) be the last
nonzero leading minor of the mutual flag pairing and
\(\delta=(\delta_h)_{1\le h\le j}\) its list of leading minors.
In this calculation \(j\ge1\). Put
\[
 e_h=\sum_{a=1}^h(s+1-a),\quad
 A_h=(\delta_h)\mathfrak d^{-2e_h},\quad
 d_{v,h}=|\delta_h|_v d^{-2e_h},\quad
 D_\delta=N A_j=\prod_{v\mid\infty}d_{v,j}^2 .
\]
Here \(v\) runs over one complex embedding above each real place
and \(|\cdot|_v\) means ordinary complex modulus. The coefficient
ideals of the two frames show that the \(A_h\) are integral.
Omit zero minors when taking their norms. For \(H_1=b^{C_Hm}\) put
\[
 B_\delta=\prod_{\substack{1\le h\le j\\\delta_h\ne0}}N A_h,\qquad
 S_\delta=\prod_v\max\left(1,\frac{\max_{1\le h\le j}d_{v,h}}
                                      {(2H_1)^{2s}}\right),\qquad N=T_u^f .
\]

Write \(e_1,\ldots,e_s,e_{s+1}=z\) for the ordered vectors
of the first augmented flag and use primes for the second. Set
\(X_0=0\), \(X_i=\langle e_1,\ldots,e_i\rangle_E\) for \(i\le s\)
and \(X=X_s\). Representatives of the same graded frame satisfy
\[
 \widetilde e_i-e_i\in X_{i-1}\quad(1\le i\le s+1),
 \qquad \widetilde e'_i-e'_i\in X'_{i-1}.
\]

These changes preserve all leading determinants by unitriangularity.
Use Lemma~\ref{lem:zero-gram-positive-rank} to choose representatives
whose first \(j\) vectors have the signed block pairing and whose
remaining representatives satisfy
\[
 I=X\cap X'=\langle e_{j+1},\ldots,e_s\rangle_E,\qquad
 e'_i=\gamma e_i\ (j<i\le s+1),\qquad\gamma\bar\gamma=1.
\]

Indeed on the common tail the two vectors have this relation
modulo the preceding flag spaces; subtracting those components
gives the equalities in the chosen representatives. For the last
vector this uses \(z'-\gamma z\in X+X'\). The active vectors
and \(z\) span a nondegenerate space. In its orthogonal choose
duals of the tail, and subtract one half of their Gram matrix
in the tail directions to make the dual span \(I^*\) isotropic.
This last change concerns rational coordinates. The split reference lattice has the following summands. On
\(X\), and on the first \(j\) vectors of \(X'\), take the
coefficient ideals \(\mathfrak d^{-(s+1-a)}\). On the dual tail
take their inverses, and on \(z\) take \(\mathcal O_E\).
Call their sum \(P_0\). The full reference lattice is
\[
 L_{\rm ref}=P_0+
 \bigl((\Lambda+X)\cap(X+X'+I^*)^\perp\bigr).
\]

Let \(U\) be the common flag unipotent group and
\(\Gamma_A=U(F)\cap\operatorname{Stab}(A)\).
For these coordinates the inverse covolume, relative to the
reference lattice, is the index ratio
\[
 \frac{[\Gamma_{L_{\rm ref}}:\Gamma_{L_{\rm ref}}\cap\Gamma_\Lambda]}
      {[\Gamma_\Lambda:\Gamma_{L_{\rm ref}}\cap\Gamma_\Lambda]} .
\]

Write \(\omega_t\) for its infimum over adapted Witt coordinates
inducing the given graded vectors. An upper estimate for this
infimum may use any one adapted system.

\begin{auxlemma}[Flag terms above one embedded datum]
\label{lem:zero-gram-frame-count}
Fix \(u,\Lambda,j,\delta\) as above and one embedded datum
\((B,\Gamma)\) of index \(x\) in the definite space of \(\Lambda\).
Two second frames define the same class if a unipotent \(n\)
of the first flag preserves \(\Lambda\) and
\(n e'_i-e''_i\in X''_{i-1}\) for every \(1\le i\le s+1\). The sum of the weights \(\omega_t\)
of distinct frame classes above this datum is at most
\begin{equation}\label{eq:zero-gram-frame-bound}
 F_\delta(x)=
 (D_\delta/x)^{k/2}8^{s+1}f^2D_\delta^{3c+6}N^2
       (B_\delta D_\delta N)^{10000(s+1)^3},
 \qquad c=1+\log_2(2f).
\end{equation}
\end{auxlemma}

\begin{proof}
\proofstep{Indices in the common orthogonal quotient.}
In \(I^\perp/I\) let \(Q\) be the image of
\(\Lambda\cap I^\perp\). It is unimodular since
\(\Lambda\cap I\) is primitive in the unimodular lattice \(\Lambda\).
Write \(P,P'\) for the images of \(\Lambda\cap X,\Lambda\cap X'\),
\(Z=P_E+P'_E\), and \(K=Q\cap Z^\perp\).
Projection along \(P_E\) and \(P'_E\) gives the lattices
\[
 K_1=(Q+P_E)\cap Z^\perp,\quad
 K_2=(Q+P'_E)\cap Z^\perp,\quad
 B_0=[Q\cap Z:P+P'],\qquad p=[K_1:K].
\]

Set \(J=\Lambda\cap I\), \(\mathfrak a_i=\mathfrak d^{-(s+1-i)}\)
and \(A_i=\Lambda\cap X_i\). The index on the first side is
\[
 b_1=[A_s:A_j+J]=\prod_{i=j+1}^{s}[\mathfrak a_i:H_i],
 \qquad H_i=\{t\in E:\exists y\in J,\ y-te_i\in X_{i-1}\}.
\]

Indeed the successive factors between \(A_j+J\) and \(A_s+J=A_s\)
are obtained by taking the leading coefficient of \(A_i\) at \(i\).
Its image is \(\mathfrak a_i\) by the frame condition, and the
image of \(J\cap X_i\) is \(H_i\). For \(i>j\), the two filtrations
of \(I\) coincide and \(e'_i=\gamma e_i\), so the corresponding
ideal on the second side is \(H'_i=\gamma^{-1}H_i\). Thus
\[
 [\mathfrak a_i:H'_i]=[\mathfrak a_i:H_i],\qquad
 |N_{E/\mathbf Q}\gamma|=1 .
\]

This uses equality of covolumes; the two ideals need not be equal.
The product gives the same index \(b_1\) on \(X'\).

Let \(\bar A_j,\bar A'_j\) be the images of the two integral
active lattices in \(I^\perp/I\). Locally choose bases adapted
to the flag filtrations. Their leading coefficient ideals are
\(\mathfrak a_1,\ldots,\mathfrak a_j\) on each side. Consequently
the absolute norm of the cross-pairing determinant, with its two
coefficient ideals, is \(D_\delta\). The full Hermitian Gram matrix
has the cross pairing and its adjoint as its two off-diagonal blocks:
\[
 \sqrt{\Delta(\bar A_j+\bar A'_j)}=D_\delta,\qquad
 [P+P':\bar A_j+\bar A'_j]=b_1^2 .
\]

Both \(K_1,K_2\) are unimodular: projection along \(P_E\)
identifies \(K_1\) isometrically with the lattice \(P^\perp/P\)
in \(Q\), and likewise for \(P'\). Locally, complete a primitive
isotropic basis to hyperbolic pairs to see unimodularity of the
quotient. Therefore \([K_1:K]=[K_2:K]=|K^\#/K|^{1/2}\).
In the unimodular lattice \(Q\), the primitive perpendicular
lattices \(K\) and \(Q\cap Z\) have discriminant groups of equal
order: projection of \(Q/(K\oplus(Q\cap Z))\) identifies it
with each discriminant group, by perfect duality. On \(P+P'\)
the square root of this order is \(D_\delta/b_1^2\), by the
off-diagonal block determinant. Passing to \(Q\cap Z\) divides
that square root by its index \(B_0\). Thus
\[
          pB_0=D_\delta/b_1^2,\qquad [K_2:K]=p .
\]

For the intersection there is also an injection
\[
 (K_1\cap K_2)/K\ \longrightarrow\ (Q\cap Z)/(P+P').
\]

Indeed, subtract the two integral lifts of the same vector along
\(P_E,P'_E\). Changing either lift adds \(P\) or \(P'\). If the
difference vanishes modulo \(P+P'\), an integral common lift
lies in \(Z^\perp\).

Inside \(Z^\perp\) take the orthogonal complement of the
distinguished negative line, and denote it by \(W\).
Let \(\Theta=K\cap W\) and \(M=M_\Gamma\). Writing \(L\)
for the target definite lattice, we have
\[
 \Theta\subset M\subset (K_1\cap K_2)\cap W,\qquad L=K_1\cap W,\qquad
 a=[M:\Theta],\qquad xa=[L:\Theta]\le p,\quad a\le B_0 .
\]

For \(\Theta\) the same integral lift supplies both entries of
the congruence defining the graph. The inequality \(xa\le p\) is
the injection \(L/\Theta\hookrightarrow K_1/K\). Similarly
\(M/\Theta\hookrightarrow(K_1\cap K_2)/K\) gives \(a\le B_0\).
To count generators, write \(y\in K_1\) as \(y=q+v\) with
\(q\in Q\), \(v\in P_E\). Sending \(y\) to \(v\bmod P\)
embeds \(K_1/K\) in \(P_E/P\), since \(Q\cap P_E=P\).
Every finite subquotient of this rank-\(j\) torsion module is
generated by at most \(j\) elements at each prime, hence globally
by the Chinese remainder theorem. Apply this to \(L/M\), a
quotient of \(L/\Theta\). Multiplication of the two index inequalities gives
\begin{equation}\label{eq:zero-gram-effective-index}
                    x(a b_1)^2\le D_\delta .
\end{equation}

\proofstep{Counting integral realizations.}
Set \(D=D_\delta,\ B_*=B_\delta,\ n=N\mathfrak c\), with
\[
 \mathfrak c=(T_u)\prod_{h:\delta_h\ne0} A_h\overline {A_h},
 \qquad n=(NB_*)^2,\qquad \ell=2s+1,\qquad e=12s+50 .
\]

In the product \(1\le h\le j\). In simultaneous echelon form
the active pairing consists of signed antidiagonal blocks.
The nonzero minors give the block ends. For a block of length
\(h\) the first entry is a root of an equation of degree \(h\)
given by the ratio of successive end determinants, up to sign.
Thus compositions, signs and roots together cost at most \(8^s\)
(\(h\le2^{h-1}\)). The primitive terminal markings in the two middle
lattices give \(K_1\cap Ez=Rz\) and \(K_2\cap Ez=Rz'\), with
\(R=\mathcal O_E\) and \(z'=\gamma z\) in the common quotient.
Use the one lattice \(K\) to define two integral ideals by
\[
 \mathfrak a=\{a\in R:az\in K\},\qquad
 \mathfrak a'=\{a\in R:az'\in K\},\qquad
                   \mathfrak a=\gamma\mathfrak a'.
\]

The map \(a\mapsto az\bmod K\) embeds \(R/\mathfrak a\) in
\(K_1/K\), and \(a\mapsto az'\bmod K\) does the same for
\(R/\mathfrak a'\) in \(K_2/K\). Thus
\(N\mathfrak a,N\mathfrak a'\le p\le D\).
For each such pair of ideals the
generators differ by norm-one units, hence by roots of unity.
Using \(\varphi(t)\ge\sqrt{t/2}\) and \([E:\mathbf Q]=2f\) gives
\[
 \#\mu(E)\le8f^2,\qquad
 \#\{\gamma\}\le 8f^2D^{2(c+2)} .
\]

Here the number of ideals of a prescribed norm \(d_0\) is at
most \(d_0^c\): at a rational prime the exponent distributions
are bounded by the ordered choices of at most \(2f\) prime factors
at each step. The first row fixes its full middle lattice \(K_1\). On the other
side the datum gives \(B=K_2\cap W\), and \(\gamma\) fixes the
marked line lattice \(Rz'\). To recover the unimodular lattice
\(K_2\) containing \(Rz'\oplus B\), its gluing must be the graph
of an anti-isometry
\[
             (Rz')^\#/Rz'\simeq R/(T_u)\longrightarrow B^\#/B .
\]

Both summands are primitive, so both graph projections are
isomorphisms by unimodularity. The image of one generator determines
the map, giving at most \(|B^\#/B|=T_u^{2f}=N^2\) choices.

Fix these choices and \(a,b_1\). A section of the quotient onto
the projective lattice \(\Theta\) lifts the definite part of the
common orthogonal lattice. Changing the Witt complement by this
section is a common unipotent transformation: the rectangular block
is the difference of the lifts and its opposite is minus its adjoint.
The quadratic correction on the common tail can be taken with
a coefficient \(t\in\mathcal O_E\) satisfying \(t+\bar t=1\).
In coordinates with \(A\) the rectangular displacement, the
quadratic block \(C\) must solve
\[
 C+C^*=-A^* A,\qquad C=-t A^* A .
\]

Such \(t\) exists because \(E/F\) is unramified at finite places
and the integral trace is locally onto, including above 2.
In the resulting orthogonal decomposition the integral definite
intersection is \(\Theta\). Write \(P_{\rm int}\) for the
intersection with its complementary rational \(\ell\)-plane.
We record the denominator calculation in this plane:
\begin{equation}\label{eq:zero-gram-lattice-sandwich}
             \mathfrak c^{e}P_0\subset P_{\rm int}
                           \subset\mathfrak c^{-e}P_0 .
\end{equation}

For the active flags, triangular elimination uses the leading
block inverses. Each inverse is the adjugate divided by the
nonzero leading determinant, with exactly the indicated powers
of the different on its rows and columns. In particular, at a
block end \(h\) each oriented denominator satisfies
\[
 \frac{\mathfrak c\,\mathfrak d^{2e_h}}{\delta_h}
       \subset\mathcal O_E,\qquad
 \frac{\mathfrak c\,\mathfrak d^{2e_h}}{\overline{\delta_h}}
       \subset\mathcal O_E .
\]

One factor of
\(\mathfrak c\) clears the denominator of that pivot or its
conjugate. Applying the row and column eliminations at each flag
step costs at most two powers. Including the negative vector
and the common-tail frame costs two more. The two active lattices
and their central reference vectors therefore have denominator
exponent \(n_0=2s+2\).

For clarity, adjoin a dual tail vector \(y\) by choosing a lift
\(\Phi y\) with the prescribed pairings. First lift the pairings on \(\mathfrak c I^*\) into the
unimodular lattice. Then project to the complementary \(\ell\)-plane.
In terms of the two projections this lift has
\[
   \Phi=\operatorname{pr}_{\ell}\circ\Psi,\qquad
     \Psi(\mathfrak c I^*_{\rm int})\subseteq\Lambda .
\]

To check the denominator of this projection, take an integral
functional on \(\Theta\), lift it to \(\Lambda\) by primitivity,
and subtract its definite component. Modulo the common tail the
difference is supported on the two active flags and \(z\).
Solving for their coefficients uses the two mutual pairing matrices
and the norm of \(z\). Successive multiplication by the ideals of
the two oriented determinants and by \((T_u)\) clears them:
\[
 \mathfrak c^3\Theta^\#\subset\Theta,\qquad
     \mathfrak c^4\Phi(I^*_{\rm int})\subset\Lambda .
\]

Here \(I^*_{\rm int}\) has the dual frame ideals. The second
inclusion spends one power on the initial lift and three on
projecting it off the definite part. Put \(w_y=\Phi y-y\), and let \(q\) be the projection onto the
common tail dual to the projection onto the chosen dual tail.
Pairing with the central vectors costs \(n_0+4\) powers. Passing
from the paired lattice to \(P_0\) costs one, since
\(\mathfrak c P_0^\#\subset P_0\) on this plane. Consequently
the denominator exponents of the following two expressions are
bounded by
\[
    w_y-qw_y:\ n_0+5,\qquad
    (w_y,w_{y'}):\ 2(n_0+5)+1 .
\]

The common-tail components pair to zero with all the \(w_y\).
Subtracting the skew part of the pairing matrix by a central
shear leaves, in the pairing with \(y'\), the error
\(t((\Phi y,\Phi y')-(w_y,w_{y'}))\).
The corrected lift differs from \(y\) by a central vector with
denominator at most \(2n_0+12\). Expressing that vector in the
integral central lattice costs a further \(n_0\) powers.
In the reverse inclusion, pair with the corrected lifts and use
\(\mathfrak c P_0^\#\subset P_0\) once more. Both exponents are at
most \(6n_0+30\le e\). All these inclusions can be tested over
the local discrete valuation rings, which proves
\eqref{eq:zero-gram-lattice-sandwich} over \(\mathcal O_E\).

It follows that the choices of the two orthogonal intersections
and of the unimodular gluing are bounded respectively by
\[
       a^{k+c},\qquad n^{2e\ell^2},\qquad ((DN)^2)^\ell .
\]
For the first bound use elementary divisors in the rank \(k\)
lattice \(M\), summing over determinant ideals of norm \(a\).
For the second choose at most \(\ell\) generators in
\(\mathfrak c^{-e}P_0/\mathfrak c^eP_0\). For the third, a
unimodular gluing is the graph of an anti-isometry of discriminant
modules. Its source has at most \(\ell\) generators. Its target
has order \((xa)^2T_u^{2f}\le(DN)^2\).

\proofstep{Relative volumes.}
Filter the common unipotent group by its action on the span of the
two flags and \(z\), then by its action on the definite complement.
A shear \(u\) fixing the definite complement pointwise acts
only on the \(\ell\)-plane. For \(\lambda\in\Lambda\), write
\(\lambda=w+y\) with \(y\) in that plane. Pairing with
\(\mathfrak c^e P_0\subset\Lambda\) gives
\[
 y\in \mathfrak c^{-e}P_0^\#\subset\mathfrak c^{-e-1}P_0,\qquad
 (u-1)P_0\subset\mathfrak c^{3e+10}P_0
 \ \Longrightarrow\
 (u-1)y\in\mathfrak c^e P_0\subset\Lambda .
\]

Self-duality gives equality \(u\Lambda=\Lambda\). This applies
directly on the central skew subgroup. On the top quotient lift
the change of the two flag spans and \(z\) to their dual-tail
coordinates while fixing the definite complement. The dual-tail
linear block is the inverse adjoint on the common tail. The
component along \(z\) is determined by its pairing divided
by \(-T_u\), and its isotropic correction is obtained with
\(t+\bar t=1\). In particular, with
\(\mathfrak n=\mathfrak c^{3e+10}\), the initial depth gives
\[
    \mathfrak c^{3e+11}\subset (T_u)\mathfrak n,\qquad
    (u z-z)/T_u\in\mathfrak n P_0 .
\]

The remaining entries of the lifted difference lie in
\(\mathfrak n P_0\): taking the inverse adjoint preserves the
identity modulo the star-stable ideal \(\mathfrak n\), and all
terms of the quadratic correction contain a changed entry.
Thus the kernel of reduction on the top block also lifts to the
integral stabilizer. Each outer quotient acts on a module of rank
at most \(\ell\), so that reduction of its matrix entries gives
\[
  \text{relative index on each outer layer}
    \ \le\ |\mathcal O_E/\mathfrak c^{3e+11}|^{\ell^2}
    = n^{(3e+11)\ell^2}.
\]

The middle quotient has one additive row for each of the
\(s-j\) common-tail vectors, given by the definite component of
the change of its dual. A row integral for both its frame ideal
and \(\Theta\) lifts to a lattice-preserving shear, by adjoining
minus its adjoint and the integral correction \(t\) above.
In row \(i\) let \(h_i\) be the index of the common-tail leading
ideal in the full frame ideal. Filtration of \(\Lambda\cap X\)
by its ordered vectors gives
\[
 \prod_{i=j+1}^s h_i
   =[\Lambda\cap X:(\Lambda\cap X_{\le j})+(\Lambda\cap I)]
   =b_1 .
\]

A change of the coefficient ideal of index \(h_i\) costs \(h_i^k\)
in a row of rank \(k\). Restricting its definite entries from
\(L\) to \(\Theta\) costs \([L:\Theta]=xa\) in each row. Thus
\[
  \text{middle relative index}\ \le\ \left(\prod_i h_i^k\right)(xa)^{s-j}
          = b_1^k(xa)^{s-j}\le b_1^k D^s .
\]

Relative indices multiply along kernels and images. Multiplying
the three lattice counts by the three relative-index bounds gives
at most
\[
 (a b_1)^k a^c(B_*DN)^{10000(s+1)^3}.
\]

Indeed \(n\le(B_*DN)^2\); the exponent after substitution is at
most \(4e\ell^2+2\ell+s+4(3e+11)\ell^2\).
There are at most \(D^2\) possible pairs \(a,b_1\).
Use \eqref{eq:zero-gram-effective-index} for the factor with
exponent \(k\), use \(a^c\le D^c\), and multiply by the block,
scalar and line counts. This gives \eqref{eq:zero-gram-frame-bound}.
\end{proof}

\begin{auxlemma}[Matched zero-Gram terms with two prescribed labels]
\label{lem:zero-gram-label-pairs}
Let \(u\) be a selected word and fix two lists \(\delta,\delta'\)
of positive active ranks \(j,j'\le s\). Take any finite set
\(\mathcal P\) of pairs of nonunit terms with \(G_t=G_{t'}=0\)
in the same finite matrix block, with those respective labels.
In particular, the two terms have the same ordered endpoint rows
\(\Lambda,\Lambda'\) and the same finite operator coordinate.
Let \(w_\Lambda\) denote the reciprocal order of the finite
row automorphism image and \(M_u=\sum_\Lambda w_\Lambda\).
Uniformly on sufficiently late rows of the marked family,
\begin{equation}\label{eq:zero-gram-weight-sum}
 \begin{gathered}
 \sum_{(t,t')\in\mathcal P}
 w_{\Lambda}w_{\Lambda'}\,\omega_t\omega_{t'}
     \le M_u E_\delta E_{\delta'},\\
 E_\delta=D_\delta^{k/2}H_C^{Cs^4}S_\delta^{16k},
       \qquad H_C=b^{Cm^2f}.
 \end{gathered}
\end{equation}
The constant \(C\) is independent of both labels and the word.
\end{auxlemma}

\begin{proof}
\proofstep{From flag pairs to embedded isometries.}
Keep the first row fixed. First form the graph in the common
quotient \(I^\perp/I\) of a single pair of flags; its integral
lattice is \(Q\) from Lemma~\ref{lem:zero-gram-frame-count}.
Set \(S=P_E+Ez\), \(S'=P'_E+Ez'\), and
\(W_t=(S+S')^\perp\) in that quotient. Write
\(L_t=(Q+P_E)\cap W_t\) and \(B_t=(Q+P'_E)\cap W_t\). For
\(y'\in B_t\), \(y\in L_t\), the finite identification is the relation
\begin{equation}\label{eq:zero-gram-congruence-graph}
 (y',y)\in\Gamma_t
 \ \Longleftrightarrow\
 \begin{cases}
 y'-l'\in P'_E,\quad y-l\in P_E\\
 \text{for some }l',l\in Q,\quad l-l'\in gQ .
 \end{cases}
\end{equation}

Indeed these are exactly the conditions on integral representatives
of the two endpoint classes modulo their respective isotropic spaces.
Using \(Q\) gives the same graph as using \(\Lambda\) before passing
to \(I^\perp/I\): integral lifts of two congruent classes differ by
an element of \(g\Lambda+(\Lambda\cap I)\), and the last summand can
be subtracted on either side.

For the original augmented flags, set \(e'_{s+1}=z'\),
\(e_{s+1}=z\). At each \(1\le i\le s+1\) and \(t\in E\), the
marking says
\[
 \forall v'\in\Lambda\cap(te'_i+X'_{i-1})\quad
 \exists v\in\Lambda\cap(te_i+X_{i-1}):
                     \quad v-v'\in g\Lambda .
\] The
condition also works in reverse: lift the desired leading coefficient
using the common frame ideal, apply the forward condition and match
the difference in the preceding space by induction. In particular,
for \(v\in Q\cap S\) the resulting \(v'\in Q\cap S'\) satisfies
\(v'\equiv v\bmod gQ\). The same assertion holds for \(Q\cap P_E\)
and \(Q\cap P'_E\).

Take \(y'\in B_t\) and choose \(l'\in Q\) with \(y'-l'\in P'_E\).
Then \(l'\perp S'\) and
\[
 \langle v,l'\rangle=\langle v-v',l'\rangle\in gR
       \qquad(v\in Q\cap S).
\]
The primitive module \(Q\cap S\) is an \(R\)-direct summand of
\(Q\) since its quotient is projective over the Dedekind domain \(R\).
Extend the integral functional \(v\mapsto\langle v,l'\rangle/g\)
to \(Q\). Unimodularity gives a representative \(m_0\in Q\).
Consequently
\[
      l=l'-gm_0\in Q\cap S^\perp .
\]
The perfect rational pairing of \(P_E\) and \(P'_E\) gives a
unique \(p\in P_E\) with \(l-p\perp P'_E\). Since \(p\perp S\)
and \(z'=\gamma z\), the vector \(y=l-p\) belongs to \(W_t\).
It gives \((y',y)\in\Gamma_t\). Interchanging the flags proves
surjectivity of the second projection as well.

If \((0,y)\in\Gamma_t\), then \(l'\in Q\cap P'_E\).
Choose \(p\in Q\cap P_E\) congruent to \(l'\). Now
\[
      (l-p)/g\in Q,\qquad y/g-(l-p)/g\in P_E,\qquad y/g\in W_t ,
\]
so \(y\in gL_t\). The reversed argument gives the other kernel
\(gB_t\); conversely \(gB_t\times gL_t\subset\Gamma_t\) follows
directly from \eqref{eq:zero-gram-congruence-graph}. Pairings of
two representatives on one side agree with the pairings of their
lifts \(l\): their differences lie in the isotropic space \(P_E\)
orthogonal to the representatives. The analogous primed lifts are
congruent modulo \(gQ\), proving preservation of the form modulo \(gR\).

For different terms the rational spaces \(W_t\) may differ.
Projection modulo the original first isotropic space canonically
identifies them with the same definite quotient
\[
        W_0=(X+Ez)^\perp/X .
\]
It is injective since \(W_t\cap P_E=0\), and surjective by the
same rational correction \(l\mapsto l-p\). Both statements concern
Hermitian spaces; projection preserves the form. Under this map \(L_t\)
becomes the fixed target lattice \(L\), the image of
\(\Lambda\cap(X+Ez)^\perp\) in \(W_0\).
Projection carries \(B_t\) to an embedded lattice \(B\) and the
two coordinates of \(\Gamma_t\) to a graph \(\Gamma\). In each middle
lattice the distinguished line is primitive of norm \(-T_u\).
Unimodular gluing as in Proposition~\ref{prop:flag-row-mass} gives
\[
 B^\#/B\simeq L^\#/L\simeq R/(T_u),\qquad
 [L:M_\Gamma]=x,\qquad L/M_\Gamma\text{ has at most }j\text{ generators}.
\]
The two index assertions follow from
Lemma~\ref{lem:zero-gram-frame-count} and its proof. Apply the
same construction separately to the other term, with index \(x'\)
and at most \(j'\) generators.

The rational isometries from the second definite row to these
two embedded lattices, say \(h_t,h_{t'}\), give
\[
   \eta=h_t h_{t'}^{-1},\qquad
   \eta B'=B,\qquad(\eta\times1)\Gamma'=\Gamma .
\]
In these equations \(h_t\) carries the fixed definite quotient
of the second row to \(B\): transport an integral representative by
the rational isometry of the term, quotient by the transported
isotropic flag, and use \(W_t\simeq W_0\). The finite operator coordinate determines the correspondence
on the two middle lattices modulo \(g\), hence on their marked
definite parts. In the ordering (source, target) this last
correspondence is given by
\[
 (v\bmod g,w\bmod g)\text{ correspond}\ \Longleftrightarrow\
 (h_t v,w)\in\Gamma
\]
by \eqref{eq:zero-gram-congruence-graph}. This proves the graph identity for \(h_th_{t'}^{-1}\).

To bound the fibre, fix a target row and one class of second
graded frames. Two realizations differ by a unipotent preserving
the target lattice, followed by a change of representatives in the
graded frame. Transporting these maps back identifies the source
lattices with their graded vectors. To identify the actual marked
vectors, split the primitive isotropic lattice \(A\) with its
hyperbolic dual as in Proposition~\ref{prop:flag-row-mass}, writing
the integral lattice as \(A\oplus H\oplus A^\vee\). In a marked
source row \(A=\bigoplus_i\mathfrak a_i e_i\). The two choices
of representatives, with these same coefficient ideals, differ on
\(A\) by an integral unitriangular automorphism \(T\); its inverse
is integral by triangular elimination. Extend it by
\[
            T\oplus 1_H\oplus (T^*)^{-1}.
\]

The remaining difference of the two marked terminal vectors lies
in \(A\), say \(a_0\). Their common image \(z_H\) is primitive
in \(H\), so there is an integral map \(C:H\to A\) with
\(Cz_H=a_0\) (split the primitive summand \(Rz_H\)). Its unitary
extension, for \(\tau+\bar\tau=1\) in \(R\), is
\[
 (a,h,y)\longmapsto
 (a+Ch-\tau CC^*y,\ h-C^*y,\ y).
\]

The two linear cross terms cancel, and the remaining quadratic
coefficient is \(1-\tau-\bar\tau=0\). Thus the source marking
can be aligned by an integral first-flag unipotent.
The two sources are the same row \(\Lambda'\) in the chosen set
of representatives. Fix one realization. For any
other, the residual source automorphism fixes the graded frame.
On \(X^\perp/X\) it fixes the negative vector and acts by an
isometry on its definite complement. Two such automorphisms with the
same induced isometry differ by a first-flag unipotent: they agree
on every graded line, the middle quotient and the dual graded lines
by unitarity. Consequently, denoting a second graded-frame class by \([F']\),
\[
 \#\{\text{terms over }\Lambda,[F']\}
        \le|\Delta_{\Lambda'}|=1/w_{\Lambda'} .
\]

Here \(\Delta_{\Lambda'}\) is the finite row automorphism image.

For a fixed pair of frame classes and a fixed triple, choice of
the first term specifies \(h_t\); the equation
\(h_{t'}=\eta^{-1}h_t\) now fixes the second term by the same
double-coset argument. Thus the weighted fibre has total source
weight at most \(w_{\Lambda'}(1/w_{\Lambda'})=1\).
The two Hilbert--Schmidt row factors are
\[
  (\sqrt{w_\Lambda w_{\Lambda'}})^2=w_\Lambda w_{\Lambda'} .
\]
The second cancels the fibre bound; the first remains for the
genus average. The relative index has the same value on
integrally equivalent graded frames, by conjugation.
Lemma~\ref{lem:zero-gram-frame-count} bounds each frame-class
sum. Proposition~\ref{prop:flag-row-mass} identifies the
normalized average on first rows with that on their definite genus.
Thus division of the left side of \eqref{eq:zero-gram-weight-sum}
by \(M_u\) leaves at most
\[
 \sum_{1\le x\le D_\delta}\sum_{1\le x'\le D_{\delta'}}
  F_\delta(x)F_{\delta'}(x')\,
         \langle|\mathcal T(L;x,x')|\rangle_{\mathcal G}.
\]

\proofstep{Absorbing the two index sums.}
Take \(C_a\ge\max(2(C_H+1),C_T,1)\), where
\(T_u\le b^{C_Tm}\). The bounds for \(m,k,f\) in Lemma~\ref{lem:rank-window}
give \(m^4/f\ll x^{-33\epsilon/4}\),
\(m^5 f/k\ll x^{-37\epsilon/16}\), and
\((\log f)^2/s\ll(\log x)^2x^{-1/2+\epsilon/16}\).
Together with the marked level choice in
Paragraph~\ref{par:window-successive-bounds} they give the late range
\[
 \begin{gathered}
 m\ge42C_a,\quad m^4\le f,\quad m^5 f\le k,\quad
 42 C_a m^2\log b\le k,\\
 \log f\le\sqrt{s},\quad 2\le d\le b,\quad (kb)^8\le g .
 \end{gathered}
\]

Set \(A=b^{C_am^2f},\ U=A S_\delta^2,\ U'=A S_{\delta'}^2,\)
and \(Z=S_\delta S_{\delta'}\).
Taking the product of the local maxima gives
\[
 D_\delta\le U,\qquad B_\delta\le U^s,\qquad N\le A,\qquad
 F_\delta(x)\le (D_\delta/x)^{k/2}U^{250000s^4}.
\]

For the last inequality substitute the first three in
\eqref{eq:zero-gram-frame-bound}; \(c\le4s\) and
\(8^{s+1},f^2\le A\). In the genus estimate
\eqref{eq:zero-gram-pair-mean}, \(Y=xx'N^4\le A^6Z^2\).
The coefficient of \(Y^7\) obeys
\[
          (32/g)^{kf}(4kf)^{4f} A^{42}\le 2^{-f}.
\]

Indeed take the \(f\)-th root, use \(b^{42C_am^2}\le e^k\),
\(f\le k,\ b\ge2\) and \(g\ge(kb)^8\); the result is at most
\((32e/(2k)^8)^k(4k^2)^4\le1/2\) for \(k\ge9\).

It follows that \(U_*,V_*\le Z^{14}/2\) in
\eqref{eq:zero-gram-pair-mean}. Since \(Z\ge1\), the truncated
sum of their powers is bounded after extracting \(Z^{14k}\)
by a geometric series of ratio \(1/2\). The two other factors obey
\[
 Y^{6s+10+4c}\le(UU')^{96s},\qquad
 2+2(20fd^2)^{2f(6+c)}d^f\le A^{183}.
\]
For the first use \(Y\le(UU')^3,\ c\le4s\); for the second
use \(f\le b^{2s},\ d\le b,\ 20\le b^5\), and \(m=2s+2\).
The genus bound contains \((xx')^{k/2}\), cancelling precisely
the denominator in the two \(F\)'s. At most \(UU'\) index
pairs occur. Hence the contribution per label is at most
\[
 D_\delta^{k/2}U^{300000s^4}S_\delta^{14k},
 \qquad 1+96s+183+250000s^4\le300000s^4 .
\]
Finally \(m^4\le f,\ m^5 f\le k,\ m\ge42\) imply
\(300000s^4\le k\). With \(C=300000C_a\) the last expression
is at most \(E_\delta\), and the same calculation applies
to \(\delta'\).
\end{proof}


\begin{auxlemma}[Summing the vanishing-Gram contribution on the diagonal]\label{lem:zero-gram-diagonal}
With \(Q=kf\log k\), the sum over the nonidentity terms at the fixed word satisfies
\[
                      \|K''_{uu}\|_{\rm HS}^2\le d_u e^{-Q}.
\]
\end{auxlemma}
\begin{proof}
Multiply the bound for each fixed label in Lemma~\ref{lem:zero-gram-fixed-label} by the corresponding arithmetic factor \(E_\delta\) of Lemma~\ref{lem:zero-gram-label-pairs}. The factors occur once for each member of the counted pair. The remaining power of \(D_\delta\) in each product is \((j+1)/2\). By Definition~\ref{def:zero-gram-label},
\[
 0<D_\delta\le (2H_1)^{4sf}S_\delta^2\le H_C S_\delta^2,\qquad
 D_\delta^{(j+1)/2}S_\delta^{-4k}
       \le H_C^{(j+1)/2}S_\delta^{-3s}.
\]


To sum the last factor over labels at fixed \(j\), put \(U_h=(2H_1)^{2s}d^{2e_h}\). Then
\[
 S_\delta^{-3s}\le\prod_{h=1}^{j}\prod_v
             \max(1,|v(\delta_h)|/U_h)^{-3}.
\]
For an integral variable \(z\), the disk packing in the proof of Lemma~\ref{lem:lattice-frequency} gives at most \(9^f(\prod_v R_v)^2\) points when every \(R_v\ge1\). 

Sum the boxes \(R_v=2^{a_v}U_h\) for \(a_v\ge0\). Their decay pays \(2^{-3\sum a_v+3f}\) and their volumes cost \(2^{2\sum a_v}U_h^{2f}\). Including zero, the sum for one coordinate is at most \(C^f U_h^{2f}\). Summing independently over all \(j\) integral coordinates costs at most \(\exp(Cm^3fL)\), since \(e_h\le s^2\). Thus the sum over labels of rank \(j\), after extracting the mass and compact dimension, has bound
\[
 B_j=\exp(Cm^6fL)(C/k)^{j(k-1-j)f}.
\]

It remains to restore the finite Gram factor in the original normalization. Write \(\Gamma_1\) for the diagonal squared norm of a finite coarse vector before division by that norm, and \(M_u=\sum_\Lambda w_\Lambda\). The product of the two finite stabilizer weights in a matching pair is \(w_\Lambda w_{\Lambda'}\) as in Lemma~\ref{lem:zero-gram-label-pairs}. 

A full marking above \(\Lambda\) is obtained by choosing a local coset of the marked subgroup in the lattice stabilizer. If \(I_\Lambda\) denotes this local coset count and \(N_\Lambda\) the number of its rational classes, orbit--stabilizer gives
\[
               I_\Lambda\le |\Delta_\Lambda|N_\Lambda .
\] The local coset number is at least \(\Gamma_1\): at the split type primes the quotient onto relative frame indices is surjective, by lifting matrices with prescribed reduction, and counting the unipotent markings in a frame gives the diagonal squared norm. Equivalently the full fibre contains at least \(\Gamma_1/|\Delta_\Lambda|\) classes. Each contributes \(\dim V_u\), by Proposition~\ref{prop:marked-column-dimension}. Hence
\[
                    \Gamma_1 M_u\dim V_u\le d_u .
\]

For a finite set of terms group first by endpoint rows and then by the finite coarse index. Orthogonality of different endpoint rows and Lemma~\ref{lem:coarse-split-gram} bound its squared norm by \(P_g\Gamma_1\) times the sum of the products of the two coefficient norms over matching pairs. Here \(P_g=\prod_{i=0}^1(1+3/g_i)^f<2\) since \(g_i\ge(kb)^8\), \(f\le k\). The preceding two label sums and their genus mean give
\[
     \|K''_{uu}\|_{\rm HS}^2\le 2d_u\Big(\sum_{j=1}^s B_j\Big)^2 .
\]


The passage to all terms follows from Lemma~\ref{lem:gram-term-sum}. On the selected window, \(s=o(k)\), \(m^6L=o(k\log k)\), and \(j(k-1-j)\ge k-2\) for \(1\le j\le s\). Consequently \(B_j\le e^{-3Q/4}\) uniformly once the row scale is sufficiently large. Since \(2s^2\le e^{Q/2}\) there, the asserted estimate follows. The argument fixes any realized target character throughout.
\end{proof}
